\documentclass[11pt]{article}
% ---------- Idioma y codificación ----------
\usepackage[utf8]{inputenc}
\usepackage[T1]{fontenc}
\usepackage[spanish]{babel}
% ---------- Página A4 ----------
\usepackage[a4paper,margin=2.5cm]{geometry}
% ---------- Bibliografía IEEE (biblatex + biber) ----------
\usepackage[backend=biber,style=ieee,sorting=none]{biblatex}
\addbibresource{referencias.bib}
% ---------- Figuras, tablas, código ----------
\usepackage{graphicx}
\usepackage{caption}
\usepackage{booktabs}
\usepackage{float}
\usepackage{xcolor}
\usepackage{listings}
\usepackage{hyperref}
\usepackage{enumitem}
\usepackage{array}
\usepackage{longtable}
\hypersetup{
    colorlinks=true,
    linkcolor=blue,
    urlcolor=blue,
    citecolor=blue
}
% Resaltado de sintaxis para SQL (Evidencia 1)
\lstdefinestyle{sqlstyle}{
    language=SQL,
    basicstyle=\ttfamily\small,
    keywordstyle=\color{blue}\bfseries,
    commentstyle=\color{gray},
    stringstyle=\color{teal},
    numbers=left,
    numberstyle=\tiny\color{gray},
    breaklines=true,
    frame=single,
    showstringspaces=false
}
\lstset{style=sqlstyle}
% Resaltado de sintaxis para Java (Evidencia 3)
\lstdefinestyle{javastyle}{
    language=Java,
    basicstyle=\ttfamily\small,
    keywordstyle=\color{blue}\bfseries,
    commentstyle=\color{gray},
    stringstyle=\color{teal},
    numbers=left,
    numberstyle=\tiny\color{gray},
    breaklines=true,
    frame=single,
    showstringspaces=false
}
\title{}
\date{}
\begin{document}
% =====================================================
% CARÁTULA
% =====================================================
    \begin{center}
        {\large\bfseries Universidad Técnica Estatal de Quevedo}\\
        Facultad de Ciencias de la Computación\\
        Ingeniería de Software --- Aplicaciones Distribuidas (ISR-701)\\[1em]
        {\Large\bfseries EV-AUT-03 --- Portafolio de Evidencias Individuales}\\[1.5em]
    \end{center}
    \noindent
    \begin{tabular}{@{}ll}
        \textbf{Estudiante:} & Jhinson Stalyn Aucatoma Celorio \\
        \textbf{Curso:} & 7mo Software ``A'' \\
        \textbf{Equipo PFC:} & AGLS \\
        \textbf{Código del instrumento:} & EV-AUT-03 \\
    \end{tabular}

    \medskip
    \noindent\textbf{Integrantes PFC:}
    \begin{itemize}[noitemsep,topsep=2pt]
          \item Jhinson Stalyn Aucatoma Celorio
          \item Jeremy Ruperto Gaibor Rodríguez
          \item Andy Paul Sánchez Pilaloa
          \item José Alejandro Lozano Morales
    \end{itemize}

    \bigskip
    \noindent\textbf{URL del repositorio del portafolio:}\\
    \url{https://github.com/JhinsonAucatoma/ev_sum_04_portafolio}

    \medskip
    \noindent\textbf{URL del repositorio del PFC:}\\
    \url{https://github.com/JoseLozanoMorales/PFC-AppsDistribuidas}

    \vspace{2em}
    \hrule
    \vspace{2em}

    \section{Evidencias de aprendizaje del periodo}

% ---------------------------------------------------------------
% EVIDENCIA 1 — CÓDIGO PROPIO
% ---------------------------------------------------------------
    \subsection{Evidencia 1 --- Código propio: criterio de fragmentación horizontal}

    \noindent\textbf{Unidad del curso:} Unidad 3. Bases de Datos Distribuidas (Tema 2: Diseño y Fragmentación de Bases de Datos Distribuidas).

    \noindent\textbf{Artefacto:}
    \begin{itemize}[noitemsep,topsep=2pt]
  \item Archivo: \texttt{sql/03\_fragmentacion\_horizontal.sql}
  \item Commit: \texttt{1640a76} (hash completo: \texttt{1640a76d196694c66c4961684013548ba3780fe6})
  \item Repositorio: \url{https://github.com/JhinsonAucatoma/practica-fragmentacion}
    \end{itemize}

    \begin{lstlisting}[caption={Fragmento de \texttt{03\_fragmentacion\_horizontal.sql}: esquema del fragmento horizontal de \texttt{pedidos}.}, label={lst:ev1-horizontal}]
CREATE TABLE pedidos_fragmento (
    pedido_id INTEGER PRIMARY KEY,
    cliente_id INTEGER NOT NULL,
    producto_id INTEGER NOT NULL,
    fecha DATE NOT NULL,
    monto NUMERIC(8,2) NOT NULL,
    sede VARCHAR(20) NOT NULL
);
    \end{lstlisting}

    \noindent\textbf{Explicación del aprendizaje.}
    Al diseñar el esquema de fragmentación horizontal de la tabla \texttt{pedidos} (Listado~\ref{lst:ev1-horizontal})
    evalué dos criterios de corte posibles: por \texttt{sede} o por \texttt{producto\_id}. Descarté \texttt{producto\_id}
    porque las consultas típicas del escenario (una cafetería con tres sedes) son del tipo ``ventas de esta sede'', no ``ventas de este producto'';
    fragmentar por producto no habría acercado los datos a quien los consulta. Entonces opté por \texttt{sede} porque ese criterio hace que cada
    fragmento resida físicamente cerca de los usuarios que más lo consultan, lo que en \emph{Principles of Distributed Database Systems} se plantea
    como uno de los objetivos centrales del diseño de fragmentación: minimizar el costo de las consultas más frecuentes acercando los datos a su patrón
    de acceso~\cite{ozsu2020principles}. Lo que aprendí es que una decisión de fragmentación no se valida solo verificando que sea correcta en abstracto
    (completitud, reconstrucción, disjunción), sino evaluando si el criterio de corte refleja cómo se va a consultar el dato en la práctica; dos esquemas
    igualmente correctos pueden tener rendimientos muy distintos según ese ajuste.

\end{document}
